% !PRINTED 561
\weberpage{561}
%<<CONT
Hierin aber darf die Reihenfolge der Summation vertauscht
werden, so dass auch

\weq{3}{\eta^{(\lambda)}\eta^{(\mu)} = \sum^{t}\sum^{s} r^{(\lambda + \mu g^{te})g^{se}}}

ist. Die nach $s$ genommene Summe

\[
\sum^{s} r^{(\lambda + \mu g^{te})g^{se}}
\]

ist aber selbst eine der Perioden, nämlich nach der Bezeichnung (1) die Periode

\[
\eta^{(\lambda + \mu g^{te})}.
\]

Es ergiebt sich also aus (2) und (3)

\weq{4}{\eta^{(\lambda)}\eta^{(\mu)} = \eta^{(\lambda+\mu)} + \eta^{(\lambda+\mu')} + \eta^{(\lambda+\mu'')} + \cdots}

Die rechte Seite dieses Ausdruckes enthält $f$ Glieder; darunter kann natürlich auch dieselbe Periode mehrmals auftreten;
auch kann darunter die uneigentliche Periode $\eta^{(0)}$ vorkommen,
die gleich der ganzen Zahl $f$ zu setzen ist. Will man die homogene Form des Ausdruckes wieder herstellen, so benutzt man
die Relation

\[
\eta + \eta_1 + \eta_2 + \cdots + \eta_{e-1} = -1,
\]

die ja nur eine andere Schreibweise der Gleichung \S.\,167, (2) ist.

Wenn man nach (4) den Ausdruck für ein Product $\eta\eta_h$
berechnet hat, dem wir oben schon die Form gegeben haben,

\weq{5}{\eta\eta_h = a_{0,h}\eta + a_{1,h}\eta_1 + \cdots + a_{e-1,h}\eta_{e-1},}

so erhält man das Product von zwei beliebigen der Perioden
durch die Substitution $(\eta, \eta_k)$,

\weq{6}{\eta_k\eta_{h+k} = a_{0,h}\eta_k + a_{1,h}\eta_{k+1} + \cdots + a_{e-1,h}\eta_{k-1}.}

Um an einem einfachen Beispiele diese Regeln zu erläutern,
nehmen wir $n = 13$. Für 13 ist 2 eine primitive Wurzel, und
wir können die im \S.\,136 gegebene Indextabelle anwenden:

\weq{7}{
\begin{array}{c|cccccccccccc}
I & 0, & 1, & 2, & 3, & 4, & 5, & 6, & 7, & 8, & 9, & 10, & 11\\
\hline
N & 1, & 2, & 4, & 8, & 3, & 6, & 12, & 11, & 9, & 5, & 10, & 7
\end{array}.
}

Wenn wir zuerst $e = 3$, $f = 4$ annehmen, so erhalten wir
eine cubische Resolvente, deren Wurzeln

\weq{8}{\begin{aligned}
\eta &= r + r_3 + r_6 + r_9\\
\eta_1 &= r_1 + r_4 + r_7 + r_{10}\\
\eta_2 &= r_2 + r_5 + r_8 + r_{11}
\end{aligned}}

sind. Aus der Tabelle (7) ergiebt sich dafür auch
%CONT>>
